% CatCrypt verified compiler: blueprint outline in leanblueprint syntax.
% Source text only. The package carries no leanblueprint build files (web.tex,
% print.tex, plastex.cfg, macros), so `leanblueprint pdf` and `leanblueprint web`
% do not run on this directory from a clone (see BUILDING.md).
% Each \lean{...}\leanok node names the fully qualified formal declaration.

\chapter{The CatCrypt verified compiler}

\section{Scope}

The nodes below are headline correctness statements of the compiler. Each is stated
for a fragment and under hypotheses, which the node text names; README.md, section
``The guarantees'', gives the full list. The front fragments are predicates over the
command language LowCT, to which hax-exported Rust is translated; no statement
concerns arbitrary Rust. The backend machine models, the
$\mathsf{decode}\circ\mathsf{encode}$ round trips and the constant-time results are
proved in Lean and are not nodes of this outline. Success of the front translation
is a hypothesis of the end-to-end statements and is established per program.

\begin{definition}[Covered control-flow fragment]
  \label{def:covered-ctl}
  \lean{CatCrypt.Crypto.SecureCompilation.Dialects.TranslateDialectDirectControl.IsCoveredCtl}
  \leanok
  The LowCT fragment of the control front: the covered straight-line arms, counted
  \texttt{RFor} loops and \texttt{RIfNz} conditionals.
\end{definition}

\begin{definition}[Scalar-covered control-flow fragment]
  \label{def:covered-sctl}
  \uses{def:covered-ctl}
  \lean{CatCrypt.Crypto.SecureCompilation.Dialects.TranslateDialectScalar.IsCoveredSCtl}
  \leanok
  The fragment of $\mathsf{IsCoveredCtl}$ on which the front's denotation is proved
  against the source semantics $\mathsf{semScalar}$: the same arms, with the side
  condition $\mathsf{srcClean}$ on the expression of a \texttt{RScalarSetExpr}
  assignment.
\end{definition}

\section{Pipeline}

The source fronts differ in the dialect they emit. The LowCT front, which receives
hax-exported Rust through $\mathsf{haxToLowCT}$, emits the SSA IR
$\mathsf{catCryptDialect}$ ($\mathsf{translateDialectControl}$), and
$\mathsf{cryptoDialectToJasminCom}$ lowers that to the word-level Jasmin dialect
$\mathsf{jasminDialect}$; its rolled-loop arms emit \texttt{Com jasminDialect}
directly. The hax front $\mathsf{compileFromHax}$ translates a hax
\texttt{ImpExpr} through THIR to \texttt{Com jasminDialect} ($\mathsf{thirToJasmin}$)
without passing through $\mathsf{catCryptDialect}$. Relational compilation from Lean
specifications emits \texttt{Com jasminDialect} (modules
\texttt{RelationalCompile} and \texttt{RelationalFront}) or
\texttt{Com fieldDialect} (module \texttt{RelationalCompileHigh}), also without
$\mathsf{catCryptDialect}$. From $\mathsf{jasminDialect}$ the backends lower to
x86-64, AArch64 and RISC-V.

\begin{theorem}[LowCT to Jasmin \texttt{Com}, end to end]
  \label{thm:rust-to-jasmin}
  \uses{def:covered-sctl}
  \lean{CatCrypt.Crypto.SecureCompilation.VIR.CompilerCorrectnessExtracted.rustToJasminCom_endToEnd_unconditional, CatCrypt.Crypto.SecureCompilation.VIR.CompilerCorrectnessExtracted.rustToJasminCom_endToEnd_topLevel_unconditional}
  \leanok
  Let $c$ be a LowCT program in $\mathsf{IsCoveredSCtl}$ on which the front
  $\mathsf{translateDialectControl}$ succeeds from the variable map $\mathit{vm}$
  with output $o$. Let $o.\mathsf{vm}$ resolve the output name to the slot $w$, let
  $\mathit{com}$ be the front's continuation closed by returning $w$
  ($o.\mathsf{cont}\ (\mathsf{Com.ret}\ w) = \mathsf{some}\ \mathit{com}$), and let
  $\mathsf{cryptoDialectToJasminCom}$ send $\mathit{com}$ to $\mathit{jcom}$. Let
  $\mathsf{semScalar}$ run $c$ from an environment consistent with the entry
  valuation $V$ to an environment holding $\mathit{val}$ at the output name, and let
  $\mathsf{RustExec}$ run $c$ from a state related to an entry valuation by
  $\mathsf{RustStateInterp}$. Then a valuation related to the final
  $\mathsf{RustExec}$ state exists, and the denotation of $\mathit{jcom}$ at the
  transport of $V$ is the transport of $\mathit{val}$. The top-level form fixes the
  empty variable map and environment and discharges both entry hypotheses.

  The value equation is stated over $\mathsf{semScalar}$, which wraps machine words
  and denotes an opaque leaf call by a default value. $\mathsf{RustExec}$ evaluates
  scalar expressions in saturating natural-number arithmetic and resolves leaf calls
  through an oracle; relating the two values on a program requires
  $\mathsf{SExpr.inRange}$ and a leaf-free body or a matching oracle, established per
  kernel.

  The theorem $\mathsf{rustToJasminCom\_endToEnd}$ in the same module, on
  $\mathsf{IsCoveredCtl}$, is conditional on its hypothesis $\mathsf{hAdequacy}$. In
  it the dialect program $\mathit{com}$ and the valuation $W$ are universally
  quantified: $\mathit{com}$ is tied to the translation output only through
  $\mathsf{hAdequacy}$, which states that $\mathit{com}$ denotes at $W$ the tracked
  output's value in every valuation related to the final state, and $W$ occurs in no
  other hypothesis. The conclusion, that $\mathit{jcom}$ denotes at the transport of
  $W$ the transport of that value, therefore takes all of its value content from
  $\mathsf{hAdequacy}$; the proof contributes the existence of the related valuation
  and the lowering equation for $\mathit{com}$.
\end{theorem}

\begin{theorem}[Robust compilation, lattice dialect to x86]
  \label{thm:robust-x86}
  \lean{CatCrypt.Crypto.SecureCompilation.VIR.RobustCompilationPipeline.robustCompilation_source_to_x86}
  \leanok
  For a program $c$ of the lattice dialect, assume leaf agreement
  $\mathsf{LeafAgrees}\ \mathit{leafSem}$, injectivity of $\mathsf{freshName}$ and of
  the parameter map, and a record $\mathit{jvc} : \mathsf{JasminVerifiedCompiler}$
  whose x86 observable, behaviour and constant-time predicate are parameters. Then the
  behaviour of $\mathsf{lowerLatticeToJasmin}\ c$ on the transported input is
  $\mathit{jvc}.\mathsf{obsOf}$ of the transported source denotation, and the lowered
  program satisfies the constant-time predicate. The two conjuncts are the fields
  $\mathsf{semRefine}$ and $\mathsf{ctPreserve}$ of $\mathit{jvc}$ applied to the
  evaluator-level and \texttt{Com}-level results; the statement contains no x86
  instruction semantics.
\end{theorem}

\begin{theorem}[Backend correctness, lattice route]
  \label{thm:lattice-nttmul-compile-correct}
  \lean{CatCrypt.Crypto.SecureCompilation.VIR.CompilerCorrectness.latticeNttMul_compile_correct}
  \leanok
  For lattice parameters $(q,n)$ and every input valuation $V$, the jasmin
  program obtained by rewriting the coefficient-form product into the NTT-domain
  sequence and lowering it across the lattice-to-jasmin conversion denotes, at the
  transported valuation, the transport of the source program's denotation. There
  is no hypothesis.

  The leaf functions of the lattice dialect take $\mathsf{ntt}$ and
  $\mathsf{invntt}$ to be the identity and $\mathsf{basemul}$ the pointwise
  product, so the theorem concerns the rewrite and the lowering, not the NTT
  arithmetic of FIPS 203. The block-executing models are
  Theorem~\ref{thm:compile-correct-basealu}.
\end{theorem}

\begin{theorem}[Backend correctness over a block-execution semantics]
  \label{thm:compile-correct-basealu}
  \lean{CatCrypt.Crypto.SecureCompilation.VIR.CompileCorrectBaseAlu.compile_correct_riscv}
  \leanok
  For a straight-line \texttt{Com} whose every operation is base-ALU
  ($\mathsf{ComBaseAlu}$) over an all-\texttt{reg 64} context, decoding the RISC-V
  model's execution of the compiled and lowered block on the encoded input state,
  where $\mathsf{exec}$ is the instruction fold $\mathsf{runRBlock64}$, reads back
  $\mathsf{Com.denote}$ at the result register, and every instruction of the block
  is base-ALU. The relation between such an instruction and an
  $\mathsf{execute\_RTYPE}$ clause copied by hand from sail-riscv is the separate
  theorem $\mathsf{ComToRiscvLsem.riscvStep64\_sail\_grounded}$.
  $\mathsf{compile\_correct\_arm}$ in the same module is the AArch64 counterpart over
  $\mathsf{runBBlock}$; no theorem relates $\mathsf{runBBlock}$ to a Sail execute
  clause.
\end{theorem}
